\clearpage
\VerdierPhysicalPageAnchor{0169}%
\begingroup\small\setlength{\parskip}{0pt}\setlength{\emergencystretch}{2.5em}
\VerdierSemanticAnchor{ch4-4-2-6-third-exact-sequence}
\[
0\longrightarrow
E\!\begin{bmatrix}n\\ f_2f_1,f_3,f_4\end{bmatrix}
\longrightarrow
E\!\begin{bmatrix}n\\ f_2,f_3,f_4\end{bmatrix}
\longrightarrow
Z\!\begin{bmatrix}n\\ f_2,f_3\end{bmatrix}
 \Big/
Z\!\begin{bmatrix}n\\ f_2f_1,f_3\end{bmatrix}
\longrightarrow0.
\]
いま $f_i$、$i\in\mathbb Z$ を、$\mathcal J$ の合成可能な射からなる
両側無限列とする。以上の議論と微分の定義から、両側無限完全列を得る：
\VerdierSemanticAnchor{ch4-4-2-6-unbounded-exact-sequence}
\[
\begin{aligned}
\cdots\longrightarrow
E\!\begin{bmatrix}n-1\\ f_{i+3},f_{i+4},f_{i+5}\end{bmatrix}
&\xrightarrow{d^{n-1}}
E\!\begin{bmatrix}n\\ f_{i+1},f_{i+2},f_{i+3}\end{bmatrix}
\longrightarrow
E\!\begin{bmatrix}n\\ f_{i+1},f_{i+3}f_{i+2},f_{i+4}\end{bmatrix}
\longrightarrow{}\\[-1pt]
&\longrightarrow
E\!\begin{bmatrix}n\\ f_{i+2},f_{i+3},f_{i+4}\end{bmatrix}
\xrightarrow{d^n}
E\!\begin{bmatrix}n+1\\ f_i,f_{i+1},f_{i+2}\end{bmatrix}
\longrightarrow\cdots.
\end{aligned}
\]
興味をもつ読者は、この完全列と古典的スペクトル系列の低次五項完全列との
関係を明らかにされたい。

\VerdierSemanticAnchor{ch4-4-2-7}
\medskip\noindent\textbf{4.2.7.}
$\mathcal A$ の任意の可換図式
\[
\begin{tikzcd}[column sep=4em,row sep=2.7em]
&Y\arrow[dl]\arrow[d,"\varphi"]\arrow[dr,"\psi"]&\\
X'\arrow[r,"\varphi'"']&X\arrow[r,"\eta"']&X'',
\end{tikzcd}
\]
で、その行が完全であるものは、$\eta$ を介して同型
（[1]、p.~316、補題 1.1）
\VerdierSemanticAnchor{ch4-4-2-7-image-isomorphism}
\[
\operatorname{Im}(\varphi)\big/\operatorname{Im}(\varphi')
\simeq\operatorname{Im}(\psi).
\]
を定める。
$\omega_3=(\mathord\bullet\xrightarrow{f_1}\mathord\bullet
\xrightarrow{f_2}\mathord\bullet\xrightarrow{f_3}\mathord\bullet)$
を $\mathcal{F}\ell_3(\mathcal J)$ の対象とする。先の考察を図式
\VerdierSemanticAnchor{ch4-4-2-7-first-diagram}
\[
\begin{tikzcd}[column sep=4.6em,row sep=2.8em]
&H^n(f_2f_1)\arrow[d]\arrow[dr]&\\
H^{n-1}(f_3)\arrow[ur,"\delta"]\arrow[r,"\delta"']&
H^n(f_2)\arrow[r]&H^n(f_3f_2),
\end{tikzcd}
\]
に適用すると、同型
\VerdierSemanticAnchor{ch4-4-2-7-1}
\[
E\!\begin{bmatrix}n\\ f_1,f_2,f_3\end{bmatrix}
\xrightarrow{\ \sim\ }
\operatorname{Im}\bigl(H^n(f_2f_1)\longrightarrow H^n(f_3f_2)\bigr).
\tag{4.2.7.1}
\]
を得る。
\endgroup
